% 04_stellar_interiors.tex
% Ch.4: Stellar impedance profiles + Neutrino MSW

\chapter{Stellar Interiors and Neutrino Oscillation}
\label{ch:stellar}

\section{Stellar Interiors as Impedance Profiles}
\label{sec:stellar_interior}

A star is a self-gravitating plasma whose electron density $n_e(r)$
increases from the corona ($\sim10^{15}\;\text{m}^{-3}$) to the core
($\sim10^{32}\;\text{m}^{-3}$).  Each radial shell has a local plasma
frequency $\omega_p(r) = \sqrt{n_e e^2/m_e\varepsilon_0}$, creating a
spherical impedance waveguide.

\begin{figure}[H]
    \centering
    \includegraphics[width=0.95\textwidth]{stellar_impedance_profile.pdf}
    \caption{Standard Solar Model mapped to an AVE impedance profile.
    \textbf{Top left:} Electron density spans 15 decades.
    \textbf{Top right:} Temperature profile (core 15.7~MK $\to$
    photosphere 5800~K $\to$ corona 1.5~MK).
    \textbf{Bottom left:} Impedance $Z(r)/Z_0$ with reflection
    coefficients at the tachocline ($\Gamma \approx 0.82$) and
    photosphere ($\Gamma \to 1$).
    \textbf{Bottom right:} Plasma frequency profile with optical
    and radio reference lines.}
    \label{fig:stellar_impedance}
\end{figure}

\subsection{Tachocline as Impedance Boundary}

The tachocline ($r \approx 0.71\,R_\odot$) separates the rigid
radiative zone from the convective envelope --- exactly analogous
to the Mohorovičić boundary in the Earth (seismic FDTD).  Its
reflection coefficient $\Gamma \approx 0.82$ traps internal
gravity waves and maintains the differential rotation profile.

\subsection{Helioseismology as Cavity Resonance}

Solar p-modes are standing acoustic waves --- transmission
line resonances of the impedance cavity, the same physics
as the protein backbone S$_{11}$ modes.

\section{Neutrino MSW: Flavor Mixing as Impedance Mode Coupling}
\label{sec:neutrino_msw}

The MSW (Mikheyev--Smirnov--Wolfenstein) effect describes neutrino
flavor oscillation in stellar matter.  In AVE, this is impedance-
dependent mode coupling between propagation channels:
\begin{itemize}
  \item $\nu_e$ mode: impedance modified by charged-current interaction
    with local electrons, $Z_e = Z_0 (1 + V_{CC}/V_0)$
  \item $\nu_\mu$ mode: impedance unchanged (no charged-current coupling)
  \item When $Z_e \approx Z_\mu$ (resonance): maximal mode conversion
\end{itemize}
The matter potential is:
\begin{resultbox}{Neutrino MSW Matter Potential}
\begin{equation}
    V_{CC} = \sqrt{2}\, G_F\, n_e
\end{equation}
\end{resultbox}
and MSW resonance occurs at the critical density:
\begin{resultbox}{MSW Resonance Critical Density}
\begin{equation}
    n_e^{res} = \frac{\Delta m^2 \cos 2\theta_{12}}{2\sqrt{2}\, G_F\, E}
\end{equation}
\end{resultbox}

\begin{figure}[H]
    \centering
    \includegraphics[width=0.95\textwidth]{neutrino_msw_profile.pdf}
    \caption{Neutrino MSW effect as impedance mode coupling.
    \textbf{Top left:} Effective mixing angle across the solar
    interior for 1--15~MeV neutrinos.
    \textbf{Top right:} Survival probability $P(\nu_e \to \nu_e)$.
    For 10~MeV neutrinos produced in the core, $P_{ee} \approx 0.31$,
    matching the SNO observation ($\sim 0.34$ for $^8$B neutrinos).
    \textbf{Bottom left:} MSW resonance density vs.\ energy.
    \textbf{Bottom right:} Mode coupling reflection coefficient
    from the universal \texttt{reflection\_coefficient()} operator.}
    \label{fig:neutrino_msw}
\end{figure}

For 10~MeV solar neutrinos, the physics engine computes
$P_{ee}(r=0) = 0.307$, consistent with the SNO measurement
of $P_{ee} = 0.340 \pm 0.023$ for $^8$B neutrinos.  The same
\texttt{reflection\_coefficient()} function that governs Pauli
exclusion, plasma cutoff, Meissner effect, seismic reflections,
and the event horizon also governs neutrino flavor conversion.
